\documentclass[10pt,a4paper]{amsart}
\usepackage{amsmath,amssymb,amsthm,graphicx}
% Retained preamble for the Stacks Project chapter morphisms.tex (Morphisms of Schemes).
% The chapter file pinned-original-source.tex carries only \input{preamble} on its first line, so no
% preamble text lives inside the chapter itself.  The Stacks Project uses one single shared preamble for
% every chapter; the snapshot stages it verbatim in sources/stacks-chow/preamble.tex.  The block below is
% copied mechanically, line for line, from that file's lines 14-18 (the Xy-pic setup) and lines 154-217
% (the native theorem declarations together with the mathematical macros).  Nothing was added, removed or
% redefined.  The wrapper class and the amsmath/amssymb/amsthm/graphicx base packages are supplied by the
% assembler, so the class/\externaldocument part of the shared preamble is not repeated here.
% For commutative diagrams we use Xy-pic
\usepackage[all]{xy}

% We use 2cell for 2-commutative diagrams.
\xyoption{2cell}
% Theorem environments.
%
\theoremstyle{plain}
\newtheorem{theorem}[subsection]{Theorem}
\newtheorem{proposition}[subsection]{Proposition}
\newtheorem{lemma}[subsection]{Lemma}

\theoremstyle{definition}
\newtheorem{definition}[subsection]{Definition}
\newtheorem{example}[subsection]{Example}
\newtheorem{exercise}[subsection]{Exercise}
\newtheorem{situation}[subsection]{Situation}

\theoremstyle{remark}
\newtheorem{remark}[subsection]{Remark}
\newtheorem{remarks}[subsection]{Remarks}

\numberwithin{equation}{subsection}

% Macros
%
\def\lim{\mathop{\mathrm{lim}}\nolimits}
\def\colim{\mathop{\mathrm{colim}}\nolimits}
\def\Spec{\mathop{\mathrm{Spec}}}
\def\Hom{\mathop{\mathrm{Hom}}\nolimits}
\def\Ext{\mathop{\mathrm{Ext}}\nolimits}
\def\SheafHom{\mathop{\mathcal{H}\!\mathit{om}}\nolimits}
\def\SheafExt{\mathop{\mathcal{E}\!\mathit{xt}}\nolimits}
\def\Sch{\mathit{Sch}}
\def\Mor{\mathop{\mathrm{Mor}}\nolimits}
\def\Ob{\mathop{\mathrm{Ob}}\nolimits}
\def\Sh{\mathop{\mathit{Sh}}\nolimits}
\def\NL{\mathop{N\!L}\nolimits}
\def\CH{\mathop{\mathrm{CH}}\nolimits}
\def\proetale{{pro\text{-}\acute{e}tale}}
\def\etale{{\acute{e}tale}}
\def\QCoh{\mathit{QCoh}}
\def\Ker{\mathop{\mathrm{Ker}}}
\def\Im{\mathop{\mathrm{Im}}}
\def\Coker{\mathop{\mathrm{Coker}}}
\def\Coim{\mathop{\mathrm{Coim}}}

% Boxtimes
%
\DeclareMathSymbol{\boxtimes}{\mathbin}{AMSa}{"02}

%
% Macros for moduli stacks/spaces
%
\def\QCohstack{\mathcal{QC}\!\mathit{oh}}
\def\Cohstack{\mathcal{C}\!\mathit{oh}}
\def\Spacesstack{\mathcal{S}\!\mathit{paces}}
\def\Quotfunctor{\mathrm{Quot}}
\def\Hilbfunctor{\mathrm{Hilb}}
\def\Curvesstack{\mathcal{C}\!\mathit{urves}}
\def\Polarizedstack{\mathcal{P}\!\mathit{olarized}}
\def\Complexesstack{\mathcal{C}\!\mathit{omplexes}}
% \Pic is the operator that assigns to X its picard group, usage \Pic(X)
% \Picardstack_{X/B} denotes the Picard stack of X over B
% \Picardfunctor_{X/B} denotes the Picard functor of X over B
\def\Pic{\mathop{\mathrm{Pic}}\nolimits}
\def\Picardstack{\mathcal{P}\!\mathit{ic}}
\def\Picardfunctor{\mathrm{Pic}}
\def\Deformationcategory{\mathcal{D}\!\mathit{ef}}

% --- appended 2026-10-12 (Chief): the project own definitions of the reference/slogan/history
% display environments, copied mechanically from the shared Stacks preamble lines 8-12
% (\usepackage{verbatim} plus the three \newenvironment lines, which the chapter preamble keeps
% verbatim from the project preamble). Nothing added, removed or redefined: these environments
% carry no mathematical content and the project itself typesets them as comments. Without these
% three lines a counted statement whose own text uses them fails with Environment slogan undefined.
\usepackage{verbatim}
\newenvironment{reference}{\comment}{\endcomment}
\newenvironment{slogan}{\comment}{\endcomment}
\newenvironment{history}{\comment}{\endcomment}
\usepackage{cleveref}
\crefname{originallemma}{original lemma}{original lemmas}
\Crefname{originallemma}{Original lemma}{Original lemmas}
\makeatletter
\newlabel{morphisms-lemma-very-ample-finite-type-over-affine}{{02NP}{source}{}{sourcepointer-1}{}}
\newlabel{morphisms-lemma-very-ample-finite-type-over-affine@cref}{{[originallemma][0][]02NP}{[1][0][]source}{}{}{}}
\makeatother
\crefname{originallemma}{original lemma}{original lemmas}
\Crefname{originallemma}{Original lemma}{Original lemmas}
\makeatletter
\newlabel{lemma-locally-finite-type-characterize}{{01T2}{source}{}{sourcepointer-2}{}}
\newlabel{lemma-locally-finite-type-characterize@cref}{{[originallemma][0][]01T2}{[1][0][]source}{}{}{}}
\makeatother
\crefname{originallemma}{original lemma}{original lemmas}
\Crefname{originallemma}{Original lemma}{Original lemmas}
\makeatletter
\newlabel{constructions-lemma-morphism-proj}{{01MY}{source}{}{sourcepointer-3}{}}
\newlabel{constructions-lemma-morphism-proj@cref}{{[originallemma][0][]01MY}{[1][0][]source}{}{}{}}
\makeatother
\title{A very ample sheaf over an affine base is pulled back from projective space}
\author{The Stacks Project Authors}
\date{stacks-project:morphisms.tex:local-snapshot-20261012; GNU FDL 1.2 or later}
\begin{document}
\maketitle
\setcounter{subsection}{1}
\noindent\textit{Editorial note.} All delivered mathematical text is the author's own, copied from the pinned Stacks Project chapter file morphisms.tex; no statement or proof step was rewritten, reordered or removed.\par\medskip
\noindent\textit{Editorial note.} The chapter file pinned-original-source.tex carries only \textbackslash{}input\{preamble\} on its first line, so no preamble text lives inside the chapter. The retained file sources/stacks-morphisms/retained-preamble.tex is a mechanical, line-for-line copy of the project's single shared preamble as staged in the snapshot. Nothing was added, removed or redefined, and the excerpt serves only as preamble, never as a span coordinate.\par\medskip
\noindent\textit{Editorial note.} Both counted spans start at the \textbackslash{}begin line of their environment, so the statement and the author's proof are delivered whole; surrounding prose of the section is deliberately outside the counted spans and is carried as standard prerequisites instead.\par\medskip
\noindent\textit{Editorial note.} The citation fields (source\_id, source\_version, authors, source\_url, license.name, source\_title) are identical to those of every other card this lane has produced from this chapter, so the batch selector treats them as one source.\par\medskip
\noindent\textit{Editorial note.} Read in full under the Lead's 12:4x rule (screening hit re-read, pb 3103B). The proof is long and complete; the author's parenthetical '(for construction theta see lemma cited above; some details omitted)' concerns one auxiliary construction inside it, not the argument of the lemma.\par\medskip
\noindent\textit{Allowed standard prerequisites (editorial).}\begin{itemize}
\item Projectivisation of a quasi-coherent module and its description as a relative Proj. Constructions, Section \textbackslash{}ref\{constructions-section-invertible-relative-proj\}, used at proof lines 9296-9305. Provides the ambient projective bundle.
\item Affine closed subschemes of the standard opens, and finite generation of the coordinate rings. Lemma \textbackslash{}ref\{lemma-locally-finite-type-characterize\}, cited at proof lines 9330-9331. Lets the local embeddings be assembled.
\end{itemize}\medskip
\ifdefined\hypertarget\hypertarget{sourcepointer-1}{}\fi
\noindent\textit{Source reference.} Original lemma 02NP: focused-seeds/AGP/sources/stacks-snapshot-20261011/tags.txt, source lines 3488--3488. The lemma's own canonical tag, quoted so that readers can resolve the label inside the delivered text; it points at the counted statement itself and carries no proof content.\par\medskip
\ifdefined\hypertarget\hypertarget{sourcepointer-2}{}\fi
\noindent\textit{Source reference.} Original lemma 01T2: focused-seeds/AGP/sources/stacks-snapshot-20261011/tags.txt, source lines 2359--2359. Navigation pointer to a statement the counted proof refers to but does not contain; nothing load-bearing is delegated to it.\par\medskip
\ifdefined\hypertarget\hypertarget{sourcepointer-3}{}\fi
\noindent\textit{Source reference.} Original lemma 01MY: focused-seeds/AGP/sources/stacks-snapshot-20261011/tags.txt, source lines 2139--2139. Navigation pointer to a statement the counted proof refers to but does not contain; nothing load-bearing is delegated to it.\par\medskip
\begin{lemma}
\label{lemma-very-ample-finite-type-over-affine}
Let $f : X \to S$ be a morphism of schemes.
Let $\mathcal{L}$ be an invertible sheaf on $X$.
Assume that
\begin{enumerate}
\item the invertible sheaf $\mathcal{L}$ is very ample on $X/S$,
\item the morphism $X \to S$ is of finite type, and
\item $S$ is affine.
\end{enumerate}
Then there exist an $n \geq 0$ and an immersion
$i : X \to \mathbf{P}^n_S$ over $S$ such that
$\mathcal{L} \cong i^*\mathcal{O}_{\mathbf{P}^n_S}(1)$.
\end{lemma}

\begin{proof}
Assume (1), (2) and (3).
Condition (3) means $S = \Spec(R)$ for some ring $R$.
Condition (1) means by definition
there exists a quasi-coherent $\mathcal{O}_S$-module
$\mathcal{E}$ and an immersion $\alpha : X \to \mathbf{P}(\mathcal{E})$
such that $\mathcal{L} = \alpha^*\mathcal{O}_{\mathbf{P}(\mathcal{E})}(1)$.
Write $\mathcal{E} = \widetilde{M}$ for some $R$-module $M$.
Thus we have
$$
\mathbf{P}(\mathcal{E}) = \text{Proj}(\text{Sym}_R(M)).
$$
Since $\alpha$ is an immersion, and since the topology of
$\text{Proj}(\text{Sym}_R(M))$ is generated by the standard
opens $D_{+}(f)$, $f \in \text{Sym}_R^d(M)$, $d \geq 1$,
we can find for each $x \in X$ an
$f \in \text{Sym}_R^d(M)$, $d \geq 1$, with $\alpha(x) \in D_{+}(f)$
such that
$$
\alpha|_{\alpha^{-1}(D_{+}(f))} : \alpha^{-1}(D_{+}(f)) \to D_{+}(f)
$$
is a closed immersion.
Condition (2) implies $X$ is quasi-compact. Hence we can find
a finite collection of elements
$f_j \in \text{Sym}_R^{d_j}(M)$, $d_j \geq 1$
such that for each $f = f_j$ the displayed map above
is a closed immersion and such that $\alpha(X) \subset \bigcup D_{+}(f_j)$.
Write $U_j = \alpha^{-1}(D_{+}(f_j))$. Note that $U_j$ is affine
as a closed subscheme of the affine scheme $D_{+}(f_j)$.
Write $U_j = \Spec(A_j)$. Condition (2) also implies that
$A_j$ is of finite type over $R$, see
Lemma \ref{lemma-locally-finite-type-characterize}.
Choose finitely many $x_{j, k} \in A_j$ which
generate $A_j$ as a $R$-algebra. Since $\alpha|_{U_j}$ is a closed
immersion we see that $x_{j, k}$ is the image of an element
$$
f_{j, k}/f_j^{e_{j, k}} \in \text{Sym}_R(M)_{(f_j)}
=
\Gamma(D_{+}(f_j), \mathcal{O}_{\text{Proj}(\text{Sym}_R(M))}).
$$
Finally, choose $n \geq 1$ and elements $y_0, \ldots, y_n \in M$ such that each
of the polynomials $f_j, f_{j, k} \in \text{Sym}_R(M)$ is a polynomial
in the elements $y_t$ with coefficients in $R$.
Consider the graded ring map
$$
\psi : R[Y_0, \ldots, Y_n] \longrightarrow \text{Sym}_R(M),
\quad Y_i \longmapsto y_i.
$$
Denote $F_j$, $F_{j, k}$ the elements of $R[Y_0, \ldots, Y_n]$ such
that $\psi(F_j) = f_j$ and $\psi(F_{j, k}) = f_{j, k}$.
By Constructions, Lemma \ref{constructions-lemma-morphism-proj}
we obtain an open subscheme
$$
U(\psi) \subset \text{Proj}(\text{Sym}_R(M))
$$
and a morphism
$r_\psi : U(\psi) \to \mathbf{P}^n_R$. This morphism
satisfies $r_\psi^{-1}(D_{+}(F_j)) = D_{+}(f_j)$, and hence we see
that $\alpha(X) \subset U(\psi)$. Moreover, it is clear
that
$$
i = r_\psi \circ \alpha : X \longrightarrow \mathbf{P}^n_R
$$
is still an immersion since
$i^\sharp(F_{j, k}/F_j^{e_{j, k}}) = x_{j, k} \in
A_j = \Gamma(U_j, \mathcal{O}_X)$
by construction. Moreover, the morphism $r_\psi$ comes
equipped with a map
$\theta : r_\psi^*\mathcal{O}_{\mathbf{P}^n_R}(1)
\to \mathcal{O}_{\text{Proj}(\text{Sym}_R(M))}(1)|_{U(\psi)}$
which is an isomorphism in this case (for construction $\theta$
see lemma cited above; some details omitted).
Since the original map $\alpha$ was assumed to have the
property that
$\mathcal{L} = \alpha^*\mathcal{O}_{\text{Proj}(\text{Sym}_R(M))}(1)$
we win.
\end{proof}

\end{document}
